Le probleme physique etudie celui de la generation de courant dans
un plasma resultant d''une faible distortion de la distribution des
electrons, il est d'usage de lineariser l'operateur de collisions
autour du noyau Maxwellien. Un discussion detaillee de cette question
est abordee dans l'article de revue de C. Karney sur les differentes
formes possibles que l'on peut utiliser pour l'operateur de collision
dans les calculs numeriques. On considerera ici trois formes, l'operateur
Maxwellien \[
C_{M}(f^{e})=C_{ee}(f^{e},f_{M}^{e})+C_{ei}(f^{e},f_{M}^{i})\]
 decrivant les collisions entre les electrons rapides et les electrons
ainsi que les ions dominants dont la distribution en impulsion est
de type Maxwellien, caracterises respectivement par les distribution
$f_{M}^{e}$ et $f_{M}^{i},$ la distribution deformee par les differentes
interactions etant ici denommee par $f^{e}.$ Pour le probleme de
la generation de courant il est important de considerer le terme de
correction d'ordre un de l'operateur de collision des electrons thermiques
sur la queue non-thermique. Dans ce cas, on parle de l'operateur de
collision linearise (ou tronque) qui est la somme de trois termes

\[
C_{M}(f^{e})=C_{ee}(f^{e},f_{M}^{e})+C_{ee}(f_{M}^{e},f^{e})+C_{ei}(f^{e},f_{M}^{i})\]


C'est l''operateur de collision generalement employe pour les simulations
plasmas. A noter qu'il introduit une legere non-linearite dans l'equation
de diffusion, ce qui necessite d'effectuer plusieurs iterations meme
en presence d'un schema numerique totalement implicite.

Enfin, il est interessant de considerer pour des comparaisons analytiques
la limite haute vitesse de l'operateur de collision \[
C_{M}^{HV}(f^{e})=C_{ee}^{HV}(f^{e},f_{M}^{e})+C_{ee}^{HV}(f_{M}^{e},f^{e})+C_{ei}^{HV}(f^{e},f_{M}^{i})\]



\subsection{Terme differentiel $C_{ee}(f,f_{M})+C_{ei}(f,f_{i})$}

L'operateur de collision decrivant les interactions binaires entre
electrons et ions est exprime sous la forme conservative \[
C_{M}(f^{e})=C_{ee}(f^{e},f_{M}^{e})+C_{ei}(f^{e},f_{M}^{i})=\frac{1}{p^{2}}\frac{\partial}{\partial p}\left[p^{2}\left(A(p)\frac{\partial f^{e}}{\partial p}+F(p)f^{e}\right)\right]+\frac{B_{t}(p)}{p^{2}}\frac{\partial}{\partial\xi}\left[(1-\xi^{2})\frac{\partial f^{e}}{\partial\xi}\right]\]


ou le terme $F(p)$ represente le ralentissement collisionnel $F(p)=F^{(e)}(p)+F^{(i)}(p)$
qui est un terme de force (convection). Normalise a la quantite $\nu_{e}p_{Te},$
le terme de collision avec les electrons vaut \[
F^{(e)}(p)=F_{1}+F_{2}\]
 ou\begin{eqnarray*}
F_{1} & = & 4\pi\int_{0}^{p}f_{M}(p^{\prime})\frac{p^{\prime}{}^{2}}{\gamma^{\prime}}\left[\frac{1}{v^{2}}+\frac{1}{p^{2}}\left(1-\frac{\gamma^{\prime}\sigma^{\prime}}{z^{\prime}}\right)\right]dp^{\prime}\\
F_{2} & = & \frac{4\pi}{v}\left(1-\frac{\sigma}{\gamma z}\right)\int_{p}^{\infty}p^{\prime}f_{M}(p^{\prime})dp^{\prime}\end{eqnarray*}
 la quantite de mouvement $p$ etant normalisee a l'impulsion thermique
$p_{Te}$, avec les definitions suivantes \[
z=p\frac{p_{Te}}{m_{e}c}\qquad\sigma=\sinh^{-1}z\qquad v=\frac{p}{\gamma}\qquad\gamma=\sqrt{1+z^{2}}\]


\begin{itemize}
\item La contribution ionique est donnee simplement par \[
F^{(i)}(p)=\frac{1}{v^{2}}Z_{i}\frac{m_{e}}{m_{i}}\]

\item \textbf{Diffusion in Momentum $p$}, described by the diffusion coefficient
$A(p)=A^{(e)}(p)+A^{(i)}(p)$, which also includes a contribution
from collisions with electrons and with ions. $A(p)$ is normalized
to the thermal diffusion coefficient $\nu_{e}p_{Te}^{2}$.

\begin{itemize}
\item \emph{Collisions with electrons} \[
A^{(e)}(p)=\frac{F_{1}+F_{2}}{v}\]

\item \emph{Collisions with ions} \[
A^{(i)}(p)=\frac{1}{v^{3}}Z_{i}\frac{m_{e}}{m_{i}}\frac{T_{i}}{T_{e}}\]

\end{itemize}
\item \textbf{Pitch-angle Scattering (Diffusion in $\xi$)}, described by
the diffusion coefficient $B_{t}(p)=B_{t}^{(e)}(p)+B_{t}^{(i)}(p)$,
normalized to $\nu_{e}p_{Te}^{2}$ with

\begin{itemize}
\item \emph{Collisions with electrons} \[
B_{t}^{(e)}(p)=B_{1}+B_{2}\]
 with \begin{eqnarray*}
B_{1} & = & 4\pi\int_{0}^{p}f_{M}(p^{\prime})\frac{p^{\prime}{}^{2}}{v}\left\{ \frac{1}{2}-\frac{p^{\prime}{}^{2}}{6p^{2}}+\frac{1}{8\gamma^{\prime}\gamma^{2}z^{2}}\left[-3\gamma^{\prime}+\sigma^{\prime}\left(\frac{3}{z^{\prime}}+2z^{\prime}\right)\right]\cdots\right.\\
 &  & \,\,\left.-\frac{1}{4z^{2}\gamma^{\prime}}\left(\gamma^{\prime}-\frac{\sigma^{\prime}}{z^{\prime}}-\frac{2}{3}\gamma^{\prime}z^{\prime}{}^{2}\right)-\frac{1}{4\gamma^{\prime}\gamma^{2}}\left(\gamma^{\prime}-\frac{\sigma^{\prime}}{z^{\prime}}\right)\right\} dp^{\prime}\\
B_{2} & = & 4\pi\int_{p}^{\infty}p^{\prime}f_{M}(p^{\prime})\frac{p^{\prime}{}^{2}}{v^{\prime}}\left\{ \frac{1}{2}-\frac{\gamma^{2}}{6\gamma^{\prime}{}^{2}}+\frac{1}{8\gamma\gamma^{\prime}{}^{2}z^{2}}\left[-3\gamma+\sigma\left(\frac{3}{z}+2z\right)\right]\cdots\right.\\
 &  & \,\,\left.-\frac{\gamma}{4z^{2}\gamma^{\prime}{}^{2}}\left(\gamma-\frac{\sigma}{z}-\frac{2}{3}\gamma z^{2}\right)-\frac{v^{\prime}{}^{2}}{4\gamma p^{2}}\left(\gamma-\frac{\sigma}{z}\right)\right\} dp^{\prime}\end{eqnarray*}

\item \emph{Collisions with ions} \[
B_{t}^{(i)}(p)=\frac{1}{2v}Z_{i}\left(1-\frac{1}{v^{2}}\frac{m_{e}}{m_{i}}\frac{T_{i}}{T_{e}}\right)\]

\end{itemize}
\end{itemize}
It is interesting to note that collisions with ions bring only a very
small contribution (ratio $m_{e}/m_{i}$) to the slowing down and
diffusion in $p$ processes. However, the contribution of ions for
pitch-angle scattering is of the same order as the contribution from
electrons.



\textbackslash{}subsubsubsection\{Expression in the Limit of a Lorentz
Gaz\}

In the limit of a Lorentz gas ($Z_{i}\gg1$), pitch-angle scattering
over ions becomes the dominant collision process and the collision
operator then reduces from () to: \[
C^{(\text{Lor.})}(f)=\frac{\nu_{e}Z}{2v}\frac{1}{p^{2}}\frac{\partial}{\partial\xi}\left[(1-\xi^{2})\frac{\partial f}{\partial\xi}\right]=\frac{\nu_{ei}}{2\gamma^{2}}\frac{\partial}{\partial\xi}\left[(1-\xi^{2})\frac{\partial f}{\partial\xi}\right]\]
 with \[
\nu_{ei}=\frac{e^{4}Zi^{2}n_{i}\ln\Lambda}{4\pi\epsilon_{0}^{2}m_{e}^{2}v^{3}}\]




\textbackslash{}subsubsubsection\{Expression as the Divergence of
a Flux and Bounce-Averaged\\
Operator\}

The differential part of the collision operator can also be expressed
as the divergence of a flux: \[
C_{ee}(f,f_{M})+C_{ei}(f,f_{i})=-\nabla\cdot\vec{S}^{C}\]
 where the expression of the divergence operator in spherical coordinates
is \[
\nabla\cdot\vec{S}=\frac{1}{p^{2}}\frac{\partial}{\partial p}(p^{2}S_{p})-\frac{1}{p}\frac{\partial}{\partial\xi}(\sqrt{1-\xi^{2}}S_{\xi})\]


Therefore, the expressions for the collisions-related fluxes are \begin{eqnarray*}
S_{p}^{C} & = & -\left[A(p)\frac{\partial f}{\partial p}+F(p)f\right]\\
S_{\xi}^{C} & = & \frac{B_{t}(p)}{p}\sqrt{1-\xi^{2}}\frac{\partial f}{\partial\xi}\end{eqnarray*}


A bounce-averaging of the collision operator is required in order
to remove the dependence upon the poloidal angle $\theta$ and reduce
the Fokker-Planck equation to a 2-D momentum space equation. This
operation is described in (), and has been derived for the collision
operator in , for the case of $f_{0}$, $g$ and $\tilde{f}$, which
gives:

\begin{itemize}
\item \emph{For $f_{0}$ or $g$} \[
\{\nabla\cdot\vec{S}^{C}\}_{f_{0},g}=\frac{1}{p^{2}}\frac{\partial}{\partial p}(p^{2}S_{p}^{C,(0)})-\frac{1}{\lambda p}\frac{\partial}{\partial\xi_{0}}(\sqrt{1-\xi_{0}^{2}}\lambda S_{\xi_{0}}^{C,(0)})\]
 where \begin{eqnarray*}
S_{p}^{C,(0)} & = & -\left[A(p)\frac{\partial f^{(0)}(p,\xi_{0})}{\partial p}+F(p)f^{(0)}(p,\xi_{0})\right]\\
S_{\xi_{0}}^{C,(0)} & = & \frac{B_{t}(p)}{p}\sqrt{1-\xi_{0}^{2}}\left(1-\frac{\Delta}{\xi_{0}^{2}}\right)\frac{\partial f^{(0)}(p,\xi_{0})}{\partial\xi_{0}}\end{eqnarray*}

\item \emph{For $\tilde{f}$} \[
\{\nabla\cdot\vec{S}^{C}\}_{\tilde{f}}=\frac{1}{p^{2}}\frac{\partial}{\partial p}(p^{2}\tilde{S}_{p}^{C,(0)})-\frac{1}{\lambda p}\frac{\partial}{\partial\xi_{0}}(\sqrt{1-\xi_{0}^{2}}\lambda\tilde{S}_{\xi_{0}}^{C,(0)})\]
 where \begin{eqnarray*}
\tilde{S}_{p}^{C,(0)} & = & -\frac{1}{\lambda(1+\epsilon)}H(|\xi_{0}|>\xi_{0T})\left[A(p)\frac{\partial\tilde{f}^{(0)}(p,\xi_{0})}{\partial p}+F(p)\tilde{f}^{(0)}(p,\xi_{0})\right]\\
\tilde{S}_{\xi_{0}}^{C,(0)} & = & \frac{1}{\lambda(1+\epsilon)}H(|\xi_{0}|>\xi_{0T})\frac{B_{t}(p)}{p}\sqrt{1-\xi_{0}^{2}}\\
 &  & \qquad\left(\frac{\partial\tilde{f}^{(0)}(p,\xi_{0})}{\partial\xi_{0}}-\frac{(\epsilon-\epsilon^{2}/2)}{(1+\epsilon)\xi_{0}}\frac{\partial}{\partial\xi_{0}}\frac{\tilde{f}^{(0)}(p,\xi_{0})}{\xi_{0}}\right)\end{eqnarray*}

\end{itemize}
The normalized bounce time $\lambda$ is defined in () and the bounce-average
coefficient $\Delta$ is defined as: \[
\Delta=\left\{ 1-\frac{1}{\Psi(\theta)}\right\} \]
 Both coefficients are described in Appendix .


\subsection{Integral term $C_{ee}(f_{M},f)$}

\label{intop}

This term is included to ensure momentum conservation in the collision
operator. For this purpose, only collisions with the first, dominant
term in the Legendre polynomials expansion need to be considered.
\[
C_{ee}(f_{M},f)=\xi I_{1}(f_{M},f_{1})\]
 with \[
I_{1}(f_{M},f_{1})=\nu_{e}f_{M}(p)\left[4\pi\frac{f_{1}(p)}{\gamma}+I_{1,1}(f_{1})+I_{1,2}(f_{1})\right]\]
 \[
f_{1}=\frac{3}{2}\int_{-1}^{1}\xi fd\xi\]
 and \begin{eqnarray*}
I_{1,1}(f_{1}) & = & 4\pi\int_{0}^{p}f_{1}(p^{\prime})\frac{p^{\prime}{}^{3}}{\gamma\gamma^{\prime}p^{2}}\left\{ \frac{1}{3}-\frac{2}{3}\gamma\gamma^{\prime}+\frac{1}{5}\gamma p^{\prime}{}^{2}+\frac{1}{p^{\prime}{}^{2}}\left(\gamma^{\prime}-\frac{\sigma^{\prime}}{z^{\prime}}\right)\right.\cdots\\
 &  & -\frac{\gamma}{z^{\prime}{}^{2}}\left(\gamma^{\prime}-\frac{\sigma^{\prime}}{z^{\prime}}-\frac{2}{3}\gamma^{\prime}z^{\prime}{}^{2}\right)+\frac{(\gamma p^{2}-5)}{6z^{\prime}}\left(-\frac{3\gamma^{\prime}\sigma^{\prime}}{z^{\prime}{}^{2}}+\frac{3}{z^{\prime}}+z^{\prime}\right)\cdots\\
 &  & +\frac{\gamma}{2z^{\prime}\beta_{th}^{2}}\left(-\frac{3\gamma^{\prime}\sigma^{\prime}}{z^{\prime}{}^{2}}+\frac{3}{z^{\prime}}+z^{\prime}-\frac{2}{5}z^{\prime}{}^{3}\right)+\frac{\gamma}{2z^{\prime}{}^{2}}\left[-3\gamma^{\prime}+\sigma^{\prime}\left(\frac{3}{z^{\prime}}+2z^{\prime}\right)\right]\cdots\\
 &  & \left.+\frac{p^{2}}{p^{\prime}{}^{2}}\left(\frac{\gamma^{\prime}\sigma^{\prime}}{z^{\prime}}-1\right)-\frac{\gamma}{12z^{\prime}\beta_{th}^{2}}\left[\gamma^{\prime}\sigma^{\prime}\left(\frac{15}{z^{\prime}{}^{2}}+6\right)-\frac{15}{z^{\prime}}-11z^{\prime}\right]\right\} dp^{\prime}\end{eqnarray*}
\begin{eqnarray*}
I_{1,2}(f_{1}) & = & 4\pi\int_{p}^{\infty}f_{1}(p^{\prime})\frac{p}{\gamma\gamma^{\prime}}\left\{ \frac{1}{3}-\frac{2}{3}\gamma\gamma^{\prime}+\frac{1}{5}\gamma^{\prime}p^{2}+\frac{1}{p^{2}}\left(\gamma-\frac{\sigma}{z}\right)\right.\cdots\\
 &  & -\frac{\gamma^{\prime}}{z^{2}}\left(\gamma-\frac{\sigma}{z}-\frac{2}{3}\gamma z^{2}\right)+\frac{(\gamma^{\prime}p^{\prime}{}^{2}-5)}{6z}\left(-\frac{3\gamma\sigma}{z^{2}}+\frac{3}{z}+z\right)\cdots\\
 &  & +\frac{\gamma^{\prime}}{2z\beta_{th}^{2}}\left(-\frac{3\gamma\sigma}{z^{2}}+\frac{3}{z}+z-\frac{2}{5}z^{3}\right)+\frac{\gamma^{\prime}}{2z^{2}}\left[-3\gamma+\sigma\left(\frac{3}{z}+2z\right)\right]\cdots\\
 &  & \left.+\frac{p^{\prime}{}^{2}}{p^{2}}\left(\frac{\gamma\sigma}{z}-1\right)-\frac{\gamma^{\prime}}{12z\beta_{th}^{2}}\left[\gamma\sigma\left(\frac{15}{z^{2}}+6\right)-\frac{15}{z}-11z\right]\right\} dp^{\prime}\end{eqnarray*}


The bounce-averaging of the integral term is derived in for $f_{0}$,
$g$ and $\tilde{f}$, which gives:

\begin{itemize}
\item \emph{For $f_{0}$ or $g$} \[
\{ C_{ee}(f_{M},f)\}_{f_{0},g}=\{\xi I_{1}(f_{M},f_{1})\}_{f_{0},g}=\frac{s^{\ast}}{\lambda}\xi_{0}I_{1}(f_{M},f_{1}^{(0)})\]
 where \[
f_{1}^{(0)}=\frac{3}{2}\int_{-1}^{1}\xi_{0}f^{(0)}d\xi_{0}\]



$s^{\ast}$ is defined as \[
s^{\ast}=H(|\xi_{0}|>\xi_{0T})\,\,\int_{-\pi}^{\pi}\frac{d\theta}{2\pi}\Psi\]
 This bounce-averaging coefficient is described in Appendix .

\item \emph{For $\tilde{f}$} \[
\{ C_{ee}(f_{M},f)\}_{\tilde{f}}=\{\xi I_{1}(f_{M},f_{1})\}_{\tilde{f}}=\frac{1}{\lambda}H(|\xi_{0}|>\xi_{0T})\xi_{0}I_{1}(f_{M},\tilde{f}_{1}^{(0)})\]
 where \[
\tilde{f}_{1}^{(0)}=\frac{3}{2}\int_{-1}^{1}H(|\xi_{0}|>\xi_{0T})\,\,\xi_{0}\tilde{s}^{\ast}\tilde{f}^{(0)}d\xi_{0}\]



$\tilde{s}^{\ast}$ is defined as \[
\tilde{s}^{\ast}=\int_{-\theta_{c}}^{\theta_{c}}\frac{d\theta}{2\pi}\frac{\xi}{\xi_{0}}\]
 and is described in Appendix.\end{itemize}

